\begin{theorem}[GaugeSpace UniformSpace sibling anchor]
\label{thm:gauge-space-uniformspace-sibling-anchor}
Every accepted $\GaugeSpaceUp$ packet exposes its finite gauge rows as a generated $\UniformSpaceUp$ entourage surface, and exposes the missing-symmetry boundary as the corresponding $\QuasiUniformSpaceUp$ directional surface. The carrier of \autoref{def:gauge-space-carrier} supplies the point row $X$, gauge rows $Q$, diagonal row $D$, refinement row $R$, optional symmetry row $S$, transport row $H$, continuation row $C$, provenance row $P$, and local naming row $N$. The entourage-generation theorem \autoref{thm:gauge-space-entourage-generation} reads $D,R,S$ as the uniform interface of \autoref{ch:concrete-instances-uniformspace-namecert}; when $S$ is absent or replaced by a directional row, the same generated gauge surface lands at the quasi-uniform boundary of \autoref{ch:concrete-instances-quasi-uniform-space-namecert}.
\end{theorem}
\begin{proof}
Project the accepted carrier through \autoref{def:gauge-space-carrier}. The gauge list $Q$ is finite, and every generated entourage request must pass through the displayed diagonal and refinement rows $D$ and $R$.

Apply \autoref{thm:gauge-space-entourage-generation}. With the symmetry row $S$ present, the generated route is exactly the $\UniformSpaceUp$ entourage surface: diagonal admission, finite refinement, symmetry transport, continuation replay, provenance, and local naming.

If the symmetry row is not present as a uniform row, the same carrier data still gives the directional diagonal-and-refinement route. That is the $\QuasiUniformSpaceUp$ boundary, because the omitted symmetry coordinate is precisely the distinction between the uniform sibling and the quasi-uniform sibling.
\end{proof}
